\subsection{Multiplicación de matrices en cadena}

La multiplicación de matrices en cadena es un ejemplo clásico de optimización de \emph{planes de ejecución}, en donde el orden en el que se ejecutan los pasos para completar una tarea influye drásticamente en su costo computacional. En este caso, el orden en el que se multiplica cada par de matrices hace que se requieran muchas más o muchas menos multiplicaciones entre escalares.

\subsubsection{Costo computacional}

\begin{tip}
  Si no recuerda las peculiaridades de la multiplicación de matrices, lea antes el \ref{apx:multiplicacion_matricial}.
\end{tip}

Primero, miremos cuántas operaciones se necesitan para multiplicar dos matrices. Para $\bmat{A}_{p\times q} \bmat{B}_{q \times r}$, el resultado será de tamaño $p\times r$, donde cada uno de los $pr$ elementos se computa mediante un producto punto entre vectores de tamaño $q$. En cada producto punto se realizan entonces $q$ multiplicaciones entre escalares, para un total de $pqr$ productos entre escalares.

Para una multiplicación en cadena de tres matrices:
\begin{itemize}
  \item Si se multiplica agrupando las matrices como $(\bmat{A}_{p\times q} \bmat{B}_{q \times r}) \bmat{C}_{r \times s}$ (que es equivalente a no agrupar, se multiplica de izquierda a derecha), el costo computacional está dado por $pqr$ del producto $(\bmat{A}\bmat{B})_{p\times r}$, al que se le suma $prs$ de la multiplicación entre ese producto y $\bmat{C}_{r \times s}$, acumulando \[pqr + prs = pr(q+s).\] Esa expresión es pequeña cuando $p$ y $r$ son pequeños comparados con $q$ y $s$.
  \item Si realiza el producto con la otra agrupación posible, $\bmat{A}_{p\times q} (\bmat{B}_{q \times r} \bmat{C}_{r \times s})$, se cuentan $qrs$ operaciones dentro del paréntesis, sumadas a $pqs$ cuando se multiplica $\bmat{A}$ por eso, para un total de \[pqs + qrs = qs(p+r)\] Esta agrupación es preferible cuando $q$ y $s$ son pequeños relativos a $p$ y $r$.
\end{itemize}
   
A partir de ese análisis es fácil fabricar ejemplos que ilustran cuán distintos son los costos computacionales con tres matrices de tamaños concretos: para el producto $\bmat{A}_{100\times 3} \bmat{B}_{3 \times 150} \bmat{C}_{150 \times 2}$, computar $\bmat{A} \bmat{B} \bmat{C}$ requiere de $100\cdot 150(3+2) = 75000$ multiplicaciones escalares, mientras que computar $\bmat{A} (\bmat{B} \bmat{C})$ requiere de tan solo $3 * 2(100 + 150) = 1500$ multiplicaciones, ¡50 veces menos!

A pesar de que el análisis anterior es útil para dimensionar por qué es importante agrupar las matrices inteligentemente, es un análisis que no es escalable: a medida que el número de matrices a multiplicar incrementa, la cantidad de agrupaciones completas crece de forma exponencial en función del número de matrices y decidir con base en fórmulas se torna inviable. Para 3 matrices solo hay 2 formas de agruparlas completamente; para 4 matrices, ya hay 5 formas: 
\begin{enumerate}
  \item $((\bmat{A} \bmat{B}) \bmat{C}) \bmat{D}$, equivalente a $\bmat{A} \bmat{B} \bmat{C} \bmat{D}$
  \item $(\bmat{A} (\bmat{B} \bmat{C})) \bmat{D}$, equivalente a $\bmat{A} (\bmat{B} \bmat{C}) \bmat{D}$
  \item $\bmat{A} ((\bmat{B} \bmat{C}) \bmat{D})$, equivalente a $\bmat{A} (\bmat{B} \bmat{C} \bmat{D})$
  \item $(\bmat{A} \bmat{B}) (\bmat{C} \bmat{D})$, equivalente a $\bmat{A} \bmat{B} (\bmat{C} \bmat{D})$
  \item $\bmat{A} (\bmat{B} (\bmat{C} \bmat{D}))$
\end{enumerate}
En general, para $n$ matrices hay $\frac{2n!}{(n+1)!n!}$ posibles agrupaciones completas. Ese cálculo corresponde a la sucesión de números de Catalán, cuya deducción es poco trivial, por lo que se desglosa en el \ref{apx:numeros_catalan}. El caso es que, enfrentados a la imposibilidad de elegir la mejor agrupación empleando solo álgebra, es hora de diseñar algoritmos.

\subsubsection{Solución recursiva}

